\documentclass[10pt,a4paper]{article}
\usepackage[utf8]{inputenc}
\usepackage[T1]{fontenc}
\usepackage[american,spanish,es-tabla]{babel}
\usepackage[margin=1.3cm,top=1.2cm,bottom=1.5cm]{geometry}
\usepackage{lmodern}
\usepackage{csquotes}
\usepackage{graphicx}
\usepackage{array}
\usepackage{tabularx}
\usepackage{booktabs}
\usepackage{float}
\usepackage{enumitem}
\usepackage{titlesec}
\usepackage{multicol}
\titleformat*{\section}{\large\bfseries}
\titlespacing*{\section}{0pt}{9pt}{3pt}
\usepackage[font=small,skip=4pt]{caption}
\usepackage[backend=biber,style=ieee,autolang=other]{biblatex}
\usepackage[hidelinks]{hyperref}
\usepackage{xurl}
\addbibresource{../referencias.bib}
\DeclareFieldFormat{sentencecase}{#1}
\AtBeginBibliography{\footnotesize}
\setlength{\parskip}{2pt}
\setlength{\parindent}{0pt}
\setlist{nosep,leftmargin=*}

\newcommand{\parte}[4]{\section{#2}{\small\textbf{Expone:} #3\quad\textbf{Tiempo:} #4\par}\smallskip\IfFileExists{secciones/#1.tex}{\input{secciones/#1}}{\textit{Falta \texttt{secciones/#1.tex}.}\par}}

\newcommand{\evidencia}[2]{\IfFileExists{../evidencias/#1}{\begin{figure}[H]\centering\includegraphics[width=0.92\linewidth,height=0.18\textheight,keepaspectratio]{../evidencias/#1}\caption{#2}\end{figure}}{\textit{Falta \texttt{evidencias/#1}.}\par}}
\newcommand{\lectura}[2]{\begin{minipage}[t]{0.32\linewidth}\centering\IfFileExists{../evidencias/#1}{\includegraphics[width=\linewidth,height=0.13\textheight,keepaspectratio]{../evidencias/#1}}{\fbox{\parbox{0.9\linewidth}{\centering\small Falta \texttt{evidencias/#1}}}}\par\vspace{2pt}{\small #2}\end{minipage}}
\newcommand{\commits}[2]{\href{https://github.com/DALBEZ-ROG/AD-EXP-1-BMR/commits/main?author=#1}{#2}}

\begin{document}

\begin{center}
  {\small Universidad Técnica Estatal de Quevedo -- Carrera de Software\par}
  {\small Aplicaciones Distribuidas (ISR-701), paralelo B\par}
  \vspace{4pt}
  {\Large\bfseries Notas de la exposición: Servicios web (SOAP y REST)\par}
  \vspace{2pt}
  {\normalsize EXP-1 (GA-SUM-02), exposición del primer corte\par}
  \vspace{4pt}
  {\small Benites Pérez Dariem Alberto \quad Mendoza Bermello Angello Agustín \quad Rivas Piloso Dayver Yahir\par}
  {\small 21 de noviembre de 2026\par}
\end{center}

\noindent{\small Tiempo planeado: 11 minutos de los 12 permitidos. Dariem Benites abre con la portada y la agenda (10 segundos).\par}

\parte{01-servicios-web}{Qué es un servicio web y a qué escala se usa}{Dariem Benites}{0:45}
\parte{02-arquitectura}{Arquitectura: SOAP y REST}{Dariem Benites}{1:05}
\parte{03-transparencias}{Transparencias de distribución}{Angello Mendoza}{0:50}
\parte{04-cap-pacelc}{CAP y PACELC}{Angello Mendoza}{0:50}
\parte{05-comunicacion}{Modelos de comunicación}{Angello Mendoza}{0:50}
\parte{06-fallos}{Tolerancia a fallos y un incidente real}{Dayver Rivas}{1:20}
\parte{07-comparacion}{Comparación y postura crítica}{Dayver Rivas}{0:50}
\parte{08-aplicacion-benites}{Aplicación en PresustentacionesAWPS}{Dariem Benites}{1:15}
\parte{09-aplicacion-mendoza}{Aplicación en SGIRI}{Angello Mendoza}{1:15}
\parte{10-aplicacion-rivas}{Aplicación en Gestión documental de titulación}{Dayver Rivas}{1:15}
\parte{11-lecciones}{Lecciones para el PFC y conclusiones}{Dariem Benites}{0:35}

\section*{Tabla de aportes}

\small
\noindent
\begin{tabularx}{\textwidth}{@{}>{\raggedright\arraybackslash}p{3.3cm} >{\raggedright\arraybackslash}p{2.2cm} >{\raggedright\arraybackslash}X@{}}
\toprule
Integrante (enlace a sus commits) & Partes que expone & Aportes \\
\midrule
\commits{DALBEZ-ROG}{Benites Pérez Dariem Alberto} & 1, 2, 8 y 11 & Repositorio, carátula, README y bitácora; base de las diapositivas y de las notas; bibliografía; secciones 1, 2, 8 y 11 con sus notas; diagrama de arquitectura; bosquejo a mano; medición REST frente a SOAP y lectura 1. En PresustentacionesAWPS, el servicio SOAP y el \texttt{404} del REST (commit \texttt{d6dfac6}). \\
\commits{Amendozab5}{Mendoza Bermello Angello Agustín} & 3, 4, 5 y 9 & Secciones 3, 4, 5 y 9 con sus notas (en la 4 cambió la cita de Brewer por la de Abadi), captura de una petición real a SGIRI, lectura 2 y bitácora. \\
\commits{drivasp}{Rivas Piloso Dayver Yahir} & 6, 7 y 10 & Secciones 6, 7 y 10 con sus notas, capturas de su API (200 con \texttt{ETag}, 304 y 201), lectura 3 y bitácora. En su proyecto agregó \texttt{ETag}, \texttt{201 Created} y \texttt{404} (commit \texttt{d816355} en \texttt{sistema-titulacion-uteq}). \\
\bottomrule
\end{tabularx}
\normalsize

\section*{Declaración de uso de inteligencia artificial}

\small
\noindent
\begin{tabularx}{\textwidth}{@{}>{\bfseries}l >{\raggedright\arraybackslash}X@{}}
  \multicolumn{2}{@{}l}{\textbf{Dariem Benites}} \\
  Herramienta: & Claude (Anthropic), modelo Claude Opus 5.5, en la aplicación Claude Code. \\
  Para qué se usó: & Me ayudo a proponer la estructura de la exposición y el reparto de tiempos, tambien verifico y rebuscó las fuentes que yo le pase para ver si tenian coherencia, armó la plantilla de Beamer y de estas notas (solo la plantilla no el contenido), la bibliografía y la carátula, y me ayudo al final de cada dia a corregir los borradores del README y de la bitácora que yo propuse. Ayudo a corregir las diapositivas y las notas de las partes generales (1 a 7 y 11), verifico que mi diagrama de arquitectura en draw.io este bien y el Ishikawa en TikZ. En PresustentacionesAWPS me ayudo a programar el servicio SOAP (contrato XSD, \emph{endpoint}, SOAP Fault y dos pruebas), cambió la respuesta del REST a \texttt{404} cuando la solicitud no existe y escribió el script de medición REST frente a SOAP. Luego me propuso opciones de preguntas críticas para los otros equipos y, antes de compilar la versión final, revisó las partes de Angello y Dayver contra el código de sus proyectos ya que eso era mucho para mi solo. \\
  Cómo se verificó: & Cada DOI se comprobó en Crossref (autores, año, título, revista, volumen, número y páginas) y cada cifra se comparó con su fuente; las páginas web y los WSDL del SRI se abrieron el 16 y 17 de noviembre de 2026. El servicio SOAP se probó con peticiones reales al backend local, y las 119 pruebas del proyecto pasaron. \\[4pt]
  \multicolumn{2}{@{}l}{\textbf{Angello Mendoza}} \\
  Herramienta: & DeepSeek Harness, modelo DeepSeek Flash. \\
  Para qué se usó: & Apoyó la corrección y el formato de mis secciones 3, 4, 5 y 9 y de mis notas, corrigió los dos diagramas de la sección 9 en TikZ que yo hice, verificó las fuentes que yo cite para ver si me servian y ordenó el repositorio. \\
  Cómo se verificó: & Los datos de SGIRI se comprobaron contra el código del proyecto, cada frase citada se sacó del PDF con su página y la captura de la API se tomó con el sistema corriendo a mano. \\[4pt]
  \multicolumn{2}{@{}l}{\textbf{Dayver Rivas}} \\
  Herramienta: & Claude (Anthropic), modelo Claude Opus 5.5. \\
  Para qué se usó: & Para corregir la sintaxis de \LaTeX{} de mis secciones, me ayudo a consultar como hacer los comandos para configurar mi entorno local y validar la estructura de los cambios de códigos de estado (\texttt{404} y \texttt{201}) y del \texttt{ETag} en la API de documentos. \\
  Cómo se verificó: & Las pruebas en Postman, la lectura 3 y su bitácora las hice yo y verifique, compilando las diapositivas y probando cada \emph{endpoint}. \\
\end{tabularx}
\normalsize

\begin{multicols}{2}
\printbibliography[title={Referencias}]
\end{multicols}

\section*{Anexo A. Primer bosquejo a mano}
\evidencia{bosquejo-a-mano.jpg}{Primer bosquejo de la arquitectura, hecho a mano por el equipo.}

\section*{Anexo B. Lectura de las fuentes}
{\small Las capturas completas están en \texttt{evidencias/}. De la lectura 2 hay dos páginas más: \texttt{lectura-2b.png} y \texttt{lectura-2c.png}.\par}
\medskip
\noindent
\lectura{lectura-1.png}{Lectura 1: Bogner et al. (2023). Dariem.}\hfill
\lectura{lectura-2.png}{Lectura 2: Abadi (2012). Angello.}\hfill
\lectura{lectura-3.png}{Lectura 3: Lercher et al. (2024). Dayver.}

\end{document}
